\clearpage
\gdef\otChapNum{0}
\pagestyle{otFront}
\phantomsection\label{demonstrations-debut}
\pdfbookmark[0]{Démonstrations de référence}{bm-demonstrations}
\otPageTitre{Démonstrations}{Des preuves à comprendre et à savoir expliquer}
\phantomsection\label{dem-second-degre}
\essSection{1}{Résoudre une équation du second degré}
\begin{demonstration}[Faire apparaître le discriminant]
Soit $f(x)=ax^2+bx+c$, avec $a\neq0$, et $\Delta=b^2-4ac$.
En développant, on vérifie l'identité
\[4af(x)=(2ax+b)^2-\Delta.\]
Ainsi $f(x)=0$ équivaut à $(2ax+b)^2=\Delta$.
Si $\Delta<0$, aucun carré réel ne convient. Si $\Delta=0$, on obtient
la seule solution $x=-b/(2a)$. Si $\Delta>0$, les deux possibilités
$2ax+b=\pm\sqrt\Delta$ donnent
\[x=\frac{-b-\sqrt\Delta}{2a}\quad\text{ou}\quad x=\frac{-b+\sqrt\Delta}{2a}.\]
\end{demonstration}
\begin{propriete}[Forme canonique et sommet]
L'identité précédente donne aussi
\[f(x)=a\left(x+\frac b{2a}\right)^2-\frac\Delta{4a}.
\]
Le sommet a donc pour abscisse $\alpha=-b/(2a)$ et pour ordonnée $f(\alpha)$.
Pour une détermination effective dans un cas simple, on complète le carré :
$x^2-6x+5=(x-3)^2-4$.
\end{propriete}
\begin{demonstration}[Somme et produit des racines]
Si les racines sont $r_1$ et $r_2$, la forme factorisée est
$f(x)=a(x-r_1)(x-r_2)$, y compris en cas de racine double.
En développant, $f(x)=ax^2-a(r_1+r_2)x+ar_1r_2$.
L'identification avec $ax^2+bx+c$ donne
$r_1+r_2=-b/a$ et $r_1r_2=c/a$.
\end{demonstration}
\begin{savoirfaire}[Expliquer une preuve]
Repérer l'hypothèse $a\neq0$, l'identité de départ et les trois cas possibles
pour le discriminant. La formule des solutions est une conséquence du calcul
sur un carré, pas une règle isolée à appliquer sans conditions.
\end{savoirfaire}
\clearpage
\otPageTitre{Démonstrations}{Suites arithmétiques, géométriques et sommes}
\phantomsection\label{dem-suites}
\essSection{2}{Termes généraux et sommes}
\begin{demonstration}[Suite arithmétique]
Passer de $u_0$ à $u_n$ nécessite $n$ ajouts de la raison $r$ :
$u_n=u_0+nr$. Entre deux indices $p$ et $n$, il y a $n-p$ ajouts,
donc $u_n=u_p+(n-p)r$.
Pour additionner $N$ termes consécutifs d'une suite arithmétique, on écrit
la somme dans les deux sens. Chaque paire « premier plus dernier »,
« deuxième plus avant-dernier », etc. a la même somme.
On obtient $2S=N(\text{premier}+\text{dernier})$.
\end{demonstration}
\begin{demonstration}[Les entiers de 1 à $n$]
Pour $n\geq1$, écrivons $S=1+2+\cdots+n$, puis $S=n+(n-1)+\cdots+1$.
En additionnant terme à terme, les $n$ colonnes valent toutes $n+1$ :
$2S=n(n+1)$, donc $S=n(n+1)/2$.
Pour $n=0$, la somme vide vaut 0 et la même formule reste valable.
\end{demonstration}
\begin{demonstration}[Suite géométrique]
Passer de $v_0$ à $v_n$ revient à multiplier $n$ fois par $q$, donc
$v_n=v_0q^n$. Pour $q=0$, cette formule donne les termes de rang $n\geq1$,
tous nuls ; le terme de rang 0 reste le terme initial donné.
Plus généralement $v_n=v_pq^{n-p}$ pour $n>p$ ; si l'on veut remonter
vers un indice inférieur, il faut $q\neq0$.
\end{demonstration}
\begin{demonstration}[Une somme géométrique]
Soit $G=1+q+\cdots+q^n$. Alors $qG=q+q^2+\cdots+q^{n+1}$.
En soustrayant, $(1-q)G=1-q^{n+1}$.
Si $q\neq1$, $G=(1-q^{n+1})/(1-q)$. Si $q=1$, il y a $n+1$ termes
égaux à 1, donc $G=n+1$. La séparation de ces deux cas évite de diviser par 0.
\end{demonstration}
\clearpage
\otPageTitre{Démonstrations}{Nombre dérivé et tangente}
\phantomsection\label{dem-derivation}
\essSection{3}{Dériver à partir d'un taux de variation}
\begin{demonstration}[La fonction carré]
Pour $f(x)=x^2$, $h\neq0$ et $a\in\R$,
\[\frac{f(a+h)-f(a)}h=\frac{2ah+h^2}h=2a+h.\]
Quand $h$ se rapproche de 0, ce taux se rapproche de $2a$.
Ainsi $f'(a)=2a$, pour tout réel $a$.
\end{demonstration}
\begin{demonstration}[La fonction inverse]
Pour $a\neq0$, $h\neq0$ et $a+h\neq0$,
\[\frac{1/(a+h)-1/a}{h}=-\frac1{a(a+h)}.\]
Quand $h$ se rapproche de 0, ce taux se rapproche de $-1/a^2$.
Ainsi la dérivée de $x\mapsto1/x$ est $-1/x^2$ sur $\R\setminus\{0\}$.
\end{demonstration}
\begin{demonstration}[Équation de la tangente]
La tangente en $A(a;f(a))$ a pour pente $f'(a)$.
Une droite de pente $m$ passant par $(a;b)$ s'écrit $y-b=m(x-a)$.
En remplaçant $m$ par $f'(a)$ et $b$ par $f(a)$, on obtient
$y=f(a)+f'(a)(x-a)$.
\end{demonstration}
\begin{exemple}
La tangente à $y=x^2$ en $a=2$ a pour équation $y=4+4(x-2)=4x-4$.
L'approximation affine associée donne $(2+h)^2\approx4+4h$ pour un petit $h$.
Le terme exact oublié est $h^2$, ce qui permet ici de mesurer l'erreur.
\end{exemple}
\clearpage
\otPageTitre{Démonstrations}{Produit et absence de dérivabilité}
\essSection{4}{Ce que signifie un nombre dérivé}
\begin{demonstration}[La fonction racine carrée en 0]
La fonction est définie pour $x\geq0$ et vaut 0 en 0. Pour $h>0$,
le taux de variation entre 0 et $h$ vaut
\[\frac{\sqrt h-0}{h}=\frac1{\sqrt h}.\]
Ces taux grandissent sans borne lorsque $h$ positif devient très petit.
Ils ne se rapprochent pas d'un nombre réel fini : la fonction racine carrée
n'est pas dérivable en 0.
\end{demonstration}
\begin{demonstration}[Dérivée d'un produit]
Soient $u$ et $v$ dérivables en $a$. Pour $h\neq0$, ajouter puis retrancher
$u(a+h)v(a)$ dans le numérateur donne
\begin{align*}
&\frac{u(a+h)v(a+h)-u(a)v(a)}h\\
&=u(a+h)\frac{v(a+h)-v(a)}h
+v(a)\frac{u(a+h)-u(a)}h.
\end{align*}
Quand $h$ se rapproche de 0, les deux taux se rapprochent de $v'(a)$ et
$u'(a)$. La dérivabilité de $u$ implique aussi que $u(a+h)$ se rapproche
de $u(a)$ : en effet $u(a+h)-u(a)$ est $h$ multiplié par un taux qui reste
proche du nombre fini $u'(a)$.
On obtient $(uv)'(a)=u(a)v'(a)+u'(a)v(a)$.
\end{demonstration}
\begin{attention}
Le raisonnement utilise la dérivabilité des deux fonctions. Il ne suffit
pas qu'elles soient définies. L'écriture de la dérivée d'un produit doit
conserver les deux contributions : l'une provenant de la variation de $u$,
l'autre de celle de $v$.
\end{attention}
\clearpage
\otPageTitre{Démonstrations}{Valeurs trigonométriques remarquables}
\phantomsection\label{dem-trigonometrie}
\essSection{5}{Construire les valeurs à partir de triangles}
\begin{demonstration}[L'angle $\pi/4$]
Dans un triangle rectangle isocèle dont les deux petits côtés mesurent 1,
Pythagore donne une hypoténuse $\sqrt2$. Les angles aigus mesurent $\pi/4$.
Le rapport « adjacent sur hypoténuse » et le rapport « opposé sur hypoténuse »
valent tous deux $1/\sqrt2=\sqrt2/2$.
Donc $\cos(\pi/4)=\sin(\pi/4)=\sqrt2/2$.
\end{demonstration}
\figureGeo[corrige]{preuve-pi4}
\begin{demonstration}[Les angles $\pi/3$ et $\pi/6$]
Partons d'un triangle équilatéral de côté 2. Une hauteur le partage en deux
triangles rectangles. La demi-base mesure 1, l'hypoténuse 2 et la hauteur
$\sqrt{2^2-1^2}=\sqrt3$.
Pour l'angle $\pi/3$, le côté adjacent mesure 1 et le côté opposé $\sqrt3$,
donc $\cos(\pi/3)=1/2$ et $\sin(\pi/3)=\sqrt3/2$.
Pour l'autre angle, $\pi/6$, les deux rôles s'échangent :
$\cos(\pi/6)=\sqrt3/2$ et $\sin(\pi/6)=1/2$.
\end{demonstration}
\begin{center}\begin{tikzpicture}[x=18mm,y=18mm]
\draw[otNavy,thick](0,0)--(2,0)--(1,1.73205)--cycle;
\draw[otGreenDark,dashed](1,1.73205)--(1,0);
\node[below]at(.5,0){$1$};\node[below]at(1.5,0){$1$};
\node[left]at(.4,.85){$2$};\node[right]at(1.6,.85){$2$};
\node[right,font=\small]at(1,.8){$\sqrt3$};
\draw[otMuted](1,.15)--(1.15,.15)--(1.15,0);
\end{tikzpicture}\end{center}
\begin{propriete}[Relier le triangle au cercle unité]
Diviser toutes les longueurs par l'hypoténuse place le point correspondant
sur le cercle unité. Ses coordonnées sont les rapports définissant cosinus
et sinus. Les symétries du cercle donnent ensuite les angles associés.
\end{propriete}
\clearpage
\otPageTitre{Démonstrations}{Produit scalaire et géométrie}
\phantomsection\label{dem-geometrie}
\essSection{6}{Deux preuves à partir des vecteurs}
\begin{demonstration}[La formule d'Al-Kashi]
Par Chasles, $\vect{BC}=\vect{AC}-\vect{AB}$.
Développer le carré de la norme donne
\[BC^2=AC^2+AB^2-2\vect{AB}\cdot\vect{AC}.
\]
Comme $\vect{AB}\cdot\vect{AC}=AB\,AC\cos\widehat{BAC}$,
on obtient $BC^2=AB^2+AC^2-2AB\,AC\cos\widehat{BAC}$.
Pour un angle droit en $A$, le cosinus est nul : on retrouve Pythagore.
\end{demonstration}
\begin{demonstration}[Le cercle de diamètre $[AB]$]
Soit $I$ le milieu de $[AB]$, avec $A\neq B$.
On a $\vect{MA}=\vect{MI}+\vect{IA}$ et
$\vect{MB}=\vect{MI}-\vect{IA}$.
Par bilinéarité et symétrie,
\[\vect{MA}\cdot\vect{MB}=MI^2-IA^2=MI^2-\frac{AB^2}{4}.
\]
Le produit est nul si et seulement si $MI=AB/2$, puisque les distances sont
positives ou nulles. C'est précisément le cercle de centre $I$ et de rayon
$AB/2$, donc de diamètre $[AB]$.
Les points $A$ et $B$ sont inclus : l'un des vecteurs du produit y est nul.
\end{demonstration}
\begin{center}\legende{Les deux configurations sont illustrées dans les repères visuels du chapitre 7.}\end{center}
\begin{savoirfaire}[Choisir un point intermédiaire]
Le milieu rend deux vecteurs opposés, ce qui supprime les termes croisés.
Le choix d'un point intermédiaire dans Chasles constitue souvent la première
étape utile d'une preuve géométrique sans coordonnées.
\end{savoirfaire}
\clearpage
\otPageTitre{Démonstrations}{Probabilités et moments d'une variable}
\phantomsection\label{dem-probabilites}
\essSection{7}{De sommes finies aux formules usuelles}
\begin{demonstration}[Probabilités totales]
Pour une partition $A_1,\ldots,A_k$, les évènements $A_i\cap B$ sont deux à deux
incompatibles et leur réunion est $B$. Ainsi $P(B)=\sum_i P(A_i\cap B)$.
Lorsque les $P(A_i)$ sont positifs, chaque terme s'écrit $P(A_i)P_{A_i}(B)$,
ce qui donne la formule des probabilités totales.
\end{demonstration}
\begin{demonstration}[Linéarité de l'espérance]
Sur un univers fini, on peut calculer $E(X)$ comme
$\sum_{\omega\in\Omega}X(\omega)P(\{\omega\})$, en regroupant ou non les issues
qui donnent la même valeur. Distribuer une somme finie donne alors
\[E(aX+bY)=\sum_{\omega\in\Omega}(aX(\omega)+bY(\omega))P(\{\omega\})
=aE(X)+bE(Y).\]
Aucune hypothèse d'indépendance n'intervient dans ce calcul.
\end{demonstration}
\begin{demonstration}[König-Huygens et transformation affine]
Avec $\mu=E(X)$, développer le carré donne
$V(X)=E(X^2-2\mu X+\mu^2)=E(X^2)-2\mu^2+\mu^2=E(X^2)-\mu^2$.
Pour $Y=aX+c$, $Y-E(Y)=a(X-E(X))$, donc
$V(Y)=a^2V(X)$ et $\sigma(Y)=|a|\sigma(X)$.
\end{demonstration}
\begin{exemple}[Une variable de Bernoulli]
Si $I$ vaut 0 ou 1 et $P(I=1)=p$, on a $I^2=I$.
Donc $E(I)=p$, $E(I^2)=p$ et $V(I)=p-p^2=p(1-p)$.
\end{exemple}
\clearpage
\pagestyle{otMain}
